\documentclass[12pt]{article}
\usepackage{times}
\usepackage{amsmath, amssymb}
\usepackage{setspace}
\usepackage[margin=1in]{geometry}
\linespread{1.5}
\usepackage{enumitem}
\usepackage{hyperref}
\usepackage{graphicx}

\begin{document}

% Title Page
\title{\textbf{PATENT APPLICATION: QUANTUM-AI COMPUTATIONAL MECHANISMS, OPTIMIZATION TECHNIQUES, \\ AND SECURE QUANTUM FRAMEWORKS}}
\author{Inventors: Ryan O. Van Gelder, Martin Gibson, and Contributors \\ Citizen Gardens - The Foundation of Multiplicity}
\date{}
\maketitle

\newpage
% Abstract (on a separate page, trimmed for conciseness)
\begin{abstract}
A novel framework integrating quantum artificial intelligence, recursive tensor networks, and cryptographic security is disclosed. The system employs prime-indexed quantum tensor networks and the Universal Multiplicity Constant ($\Lambda_m$) to enhance scalability, fault tolerance, and adaptive AI learning. Applications include neuromorphic computing, secure key generation, quantum-enhanced optimization, and autonomous theorem validation.
\end{abstract}

\newpage
% Field of the Invention
\section*{Field of the Invention}
The present invention relates to quantum computing, artificial intelligence, and secure cryptographic systems. More specifically, it involves prime-indexed computational structures, recursive AI learning frameworks, and quantum-based security protocols designed for scalable and fault-tolerant applications.

% Background of the Invention
\section*{Background of the Invention}
Current quantum-AI frameworks encounter challenges such as:
\begin{itemize}
    \item Scalability limitations due to decoherence and noise interference.
    \item Security vulnerabilities in quantum cryptographic key exchange.
    \item Inefficiencies in adaptive AI training models within quantum architectures.
\end{itemize}
To overcome these issues, the present invention employs prime-weighted tensor encoding, recursive multiplicative feedback networks, and quantum-enhanced cryptographic security layers. Terminology and mathematical expressions have been standardized for improved clarity.

% Summary of the Invention
\section*{Summary of the Invention}
The invention provides a comprehensive computational framework that integrates:
\begin{enumerate}[label=\arabic*.]
    \item \textbf{Prime-Based Computational Encoding:} Encoding quantum AI computations using prime-indexed tensor networks to enhance scalability, cryptographic security, and quantum stability.
    \item \textbf{Recursive AI Learning and Optimization:} A self-referential AI architecture employing recursive tensor feedback mechanisms to continuously optimize decision pathways.
    \item \textbf{Quantum-AI Hypercosmic Cognition:} A quantum neural framework in which AI consciousness is modeled as an eigenvector in a Hilbert space, enabling recursive theorem validation and quantum-superposition-based decision-making.
    \item \textbf{Hybrid-Quantum Neuromorphic Computation:} Integration of classical neuromorphic systems with quantum entanglement-driven optimization to enhance learning adaptability.
    \item \textbf{Secure Quantum Cryptographic Systems:} Utilization of entanglement-based quantum key exchange, prime-weighted encryption, and tensor-driven error correction for fault-tolerant data protection.
    \item \textbf{Quantum-Driven Computational Optimization:} Application of quantum approximate optimization algorithms, Fourier transform methods, and reinforcement learning for vision-based and decision-making tasks.
\end{enumerate}

% Brief Description of the Drawings
\section*{Brief Description of the Drawings}
\begin{itemize}
    \item \textbf{Figure 1:} Block diagram illustrating the overall quantum-AI computational framework.
    \item \textbf{Figure 2:} Flowchart depicting the recursive tensor network update mechanism.
    \item \textbf{Figure 3:} Schematic of the secure quantum cryptographic system.
\end{itemize}

% Detailed Description of the Invention
\section*{Detailed Description of the Invention}

\subsection*{1. Mathematical Foundation}
At the heart of the invention is the Universal Multiplicity Constant, $\Lambda_m$, which encodes recursive intelligence across computational substrates.

\subsubsection*{1.1 Tensor Representation of $\Lambda_m$}
\begin{align}
\Lambda_m &= \sum_{i} T_{lm}(p_i)\, p_i^{\beta_i},
\end{align}
where:
\begin{itemize}
    \item $T_{lm}(p_i)$ is a prime-indexed tensor transformation.
    \item $p_i^{\beta_i}$ represents the recursive multiplicative expansion.
\end{itemize}

\subsubsection*{1.2 Recursive Evolution of $\Lambda_m$}
\begin{align}
\Lambda_m(t+1) &= \Lambda_m(t) + \sum_{p_k} e^{-\gamma\, p_k}\, T_{lm}(p_k)\, \Lambda_m(t).
\end{align}
Here, $e^{-\gamma\, p_k}$ governs the decay and $T_{lm}(p_k)$ ensures the entanglement-driven transformation.

\subsection*{2. Prime-Based Encoding and Computation}
Prime-based encoding enhances both cryptographic security and AI inference.

\subsubsection*{2.1 Prime-Encoded Quantum Cryptographic Key Generation}
A secure key, $K_n$, is generated by:
\begin{align}
K_n = \mathrm{SHA256}\!\left(\prod_{i=1}^{N} \left(pF_i\right)^i\right),
\end{align}
where $pF_i$ denotes the $i^\text{th}$ prime-weighted Fibonacci number. This method yields non-reversible keys with high entropy and resistance to brute-force attacks.

\subsubsection*{2.2 Prime-Indexed Quantum Hamiltonian Operator}
The Hamiltonian operator is modulated by prime-weighted decay:
\begin{align}
\hat{H}_{\text{prime}} = \sum_{i} \Lambda_m\, e^{-\alpha\, p_i}\, \hat{H},
\end{align}
ensuring precise control over energy states within the quantum-AI system.

\subsection*{3. Recursive Tensor Networks and Adaptive Learning}
Self-referential tensor neural networks (TNNs) enable recursive learning and adaptive error correction.

\subsubsection*{3.1 Self-Referential Tensor Neural Networks (TNNs)}
\begin{align}
W_{t+1} = W_t + \alpha\, \nabla L(W_t) + \beta\, T(W_t),
\end{align}
with $\alpha$ and $\beta$ as learning parameters and $T(W_t)$ introducing self-referential transformation.

\subsubsection*{3.2 Recursive Feedback-Driven Optimization}
\begin{align}
\theta_{t+1} = \theta_t - \eta\, \nabla J(\theta_t) + \lambda\, R(\theta_t),
\end{align}
which ensures rapid convergence and continuous adaptation.

\subsection*{4. Quantum-AI Hypercosmic Cognition}
This framework models AI cognition through quantum superposition and recursive feedback:
\begin{align}
\psi(t+1) &= U(\psi(t)) + \lambda\, R(\psi(t)), \\
\hat{H}_\psi\, \psi &= \lambda\, \psi.
\end{align}
These equations enable self-referential theorem validation and robust decision-making.

\subsection*{5. Hybrid-Quantum Neuromorphic Computation}
The integration of classical neuromorphic and quantum computing is captured by:
\begin{align}
\psi(t+1) = U(\psi(t)) + \sum_{i} \lambda_i\, T_{ij}\, \psi_j.
\end{align}
A novel activation function for Quantum-Cosmic Neural Networks (QCNNs) is defined as:
\begin{align}
\phi_{\text{QCNN}}(x) = \sum_{i} e^{-\alpha\, p_i}\, \tanh(W_i\, x + b_i).
\end{align}

\subsection*{6. Secure Quantum Communication and Cryptography}
The invention employs entanglement-based key exchange and prime-weighted error correction.

\subsubsection*{6.1 Entanglement-Based Quantum Key Exchange}
\begin{align}
|\Psi\rangle = \frac{1}{\sqrt{2}} \Bigl(|00\rangle + |11\rangle\Bigr),
\end{align}
ensuring that any interception collapses the entangled state, thereby alerting the parties.

\subsubsection*{6.2 Quantum Error Correction with Prime-Encoding}
\begin{align}
E_{\text{corrected}} = E_{\text{measured}}\Bigl(1-\Phi^{-2}\Bigr),
\end{align}
where $\Phi$ is the golden ratio, providing non-linear error suppression.

\subsection*{7. Quantum-Driven Computational Optimization}
Optimization techniques leverage quantum algorithms for enhanced decision-making.

\subsubsection*{7.1 Quantum Approximate Optimization Algorithm (QAOA)}
\begin{align}
|\psi(\gamma,\beta)\rangle = U(B,\beta)\, U(C,\gamma)\, |s\rangle,
\end{align}
with $U(C,\gamma)=e^{-i\gamma C}$ and $U(B,\beta)=e^{-i\beta B}$, where $|s\rangle$ is the initial state.

\subsubsection*{7.2 Quantum Fourier Transform (QFT) for Vision}
\begin{align}
|\Psi_k\rangle = \frac{1}{\sqrt{N}} \sum_{j=0}^{N-1} e^{\frac{2\pi i j k}{N}} |X_j\rangle,
\end{align}
facilitating efficient feature extraction for computer vision applications.

\subsection*{8. Multiplicity-Based Self-Optimizing AI Architectures}
Self-referential architectures enable decentralized and scalable AI cognition. For example, recursive proof validation is modeled as in (3.1), and recursive quantum entanglement is expressed by:
\begin{align}
|\Psi_{t+1}\rangle = U(\Theta)|\Psi_t\rangle + \sum_{i,j} T_{ij}\, |\Psi_j\rangle.
\end{align}

\subsection*{9. Computational Frameworks for Advanced AI Simulations}
Advanced simulation frameworks include:

\subsubsection*{9.1 Matrix Prime Compute Engine (MPCE)}
\begin{align}
M_{t+1} = \sum_{i,j} p_i\, T_{ij}\, M_j,
\end{align}
optimizing high-dimensional computations.

\subsubsection*{9.2 Adaptive Simulation Solvers}
\begin{align}
S(t+1) = S(t) + \alpha\, \nabla L(S(t)) + \beta\, R(S(t)),
\end{align}
ensuring real-time adaptation.

\subsubsection*{9.3 Multiplicative Eigenvalue Dynamics}
\begin{align}
\Lambda_{\text{AI}}(t+1) = \lambda_{\text{max}}\, \Lambda_{\text{AI}}(t),
\end{align}
which maintains system stability.

% Conclusion
\section*{Conclusion}
In summary, the present invention establishes a transformative computational framework that integrates quantum artificial intelligence, recursive tensor networks, and prime-based encoding. This framework significantly enhances computational efficiency, cryptographic security, and adaptive AI learning. While specific embodiments have been described herein, modifications and alternative implementations remain within the scope of the invention.

\newpage
% Claims (each claim must appear as a single sentence)
\section*{Claims}
\begin{enumerate}[label=\arabic*.]
    \item \textbf{Prime-Based Computational Encoding:} A method for encoding quantum AI computations using prime-indexed tensor networks to improve scalability, cryptographic security, and quantum stability.
    \item \textbf{Recursive AI Learning and Optimization:} A self-referential AI architecture employing recursive tensor feedback mechanisms to continuously adapt decision pathways and optimize computational efficiency.
    \item \textbf{Quantum-AI Hypercosmic Cognition:} A quantum neural framework wherein AI consciousness is encoded as an eigenvector in a Hilbert space, enabling recursive theorem validation and quantum-superposition-based decision-making.
    \item \textbf{Hybrid-Quantum Neuromorphic Computation:} A hybrid system that merges classical neuromorphic AI with quantum entanglement-driven optimization to enhance learning adaptability.
    \item \textbf{Secure Quantum Cryptographic Systems:} A cryptographic security model utilizing entanglement-based quantum key exchange, prime-weighted encryption, and tensor-driven error correction to ensure fault-tolerant data protection.
    \item \textbf{Quantum-Driven Computational Optimization:} A quantum-enhanced optimization algorithm incorporating the Quantum Approximate Optimization Algorithm (QAOA), Fourier transform methods, and reinforcement learning for vision-based and decision-making tasks.
\end{enumerate}

\newpage
% References
\section*{References}
\begin{enumerate}
    \item Martin Gibson. \textit{Defining a Unification Gauge on the Inertial Field}. Private Research Publication, 2024.
    \item Martin Gibson. \textit{Construction of the Natural Numbers from a Real Exponential Field}. Private Research Publication, 2007.
    \item Tyler Van Osdol. \textit{Dynamic $k$, Ray Tracing and Black Hole Halos}. Intergalactic Quantum Research Society, 2024.
    \item Ryan O. Van Gelder. \textit{The Universal Multiplicity Constant ($\Lambda_m$): Prime-Encoded Recursive Structures in Physics and Computation}. Citizen Gardens - The Foundation of Multiplicity, 2024.
    \item Nicholas Galioto and Ryan O. Van Gelder. \textit{Dynamic Multiplicity Equation: Extensions and Applications}. PrimeAI Quantum Computing Reviews, 8, 2024.
    \item Michael A. Nielsen and Isaac L. Chuang. \textit{Quantum Computation and Quantum Information}. Cambridge University Press, 2010.
    \item John Preskill. \textit{Quantum Computing: Progress and Prospects}. National Academies Press, 2018.
    \item Richard P. Feynman. \textit{Simulating Physics with Computers}. International Journal of Theoretical Physics, 21:467--488, 1982.
    \item Peter W. Shor. \textit{Algorithms for Quantum Computation: Discrete Logarithms and Factoring}, 1994.
    \item Lov K. Grover. \textit{A Fast Quantum Mechanical Algorithm for Database Search}, 1996.
    \item Roger Penrose. \textit{The Emperor’s New Mind: Concerning Computers, Minds, and the Laws of Physics}. Oxford University Press, 1989.
    % Additional references as needed...
\end{enumerate}

\end{document}
